\section{Received Prayers, Collects, Antiphons, and Hymns}\label{app:prayers}

Every prayer below reproduces an identified received text without modern
alteration or conflation. Full publication details for every Latin and English
source stand in Appendix~\ref{app:references}.

\subsection{Collects of Christmastide}

All Collects of Christmastide are addressed to God the Father; their Latin
wording is checked directly against the 1962 typical edition of the
\textit{Missale Romanum}.

\begin{massbox}{Christmas Day: First Mass at Midnight (25 December)}
\textbf{Latin} (1962 Roman Missal). Deus, qui hanc sacratissimam noctem veri
luminis fecisti illustratione clarescere: da, quaesumus; ut, cujus lucis
mysteria in terra cognovimus, ejus quoque gaudiis in caelo perfruamur: Qui
tecum vivit et regnat in unitate.

\textbf{English} (Lasance, 1945). O God, Who hast brightened this most holy
night with the shining of the true light, grant, we beseech Thee, that we may
enjoy in heaven the delights of Him whose mystical light we have known on earth.
Who with Thee.
\end{massbox}

\begin{massbox}{Christmas Day: Second Mass at Dawn (25 December)}
\textbf{Latin} (1962 Roman Missal). Da nobis, quaesumus, omnipotens Deus: ut,
qui nova incarnati Verbi tui luce perfundimur; hoc in nostro resplendeat opere,
quod per fidem fulget in mente. Per eundem Dominum nostrum Jesum Christum,
Filium tuum: Qui tecum vivit et regnat in unitate.

\textbf{English} (Lasance, 1945). Grant, we beseech Thee, O almighty God, that
we who are filled with the new light of Thine Incarnate Word, may show forth in
our deeds that which by faith shineth in our minds. Through the same.
\end{massbox}

\begin{massbox}{Commemoration of St.~Anastasia (Dawn Mass, 25 December)}
\textbf{Latin} (1962 Roman Missal). Da, quaesumus, omnipotens Deus: ut, qui
beatae Anastasiae Martyris tuae solemnia colimus; ejus apud te patrocinia
sentiamus. Per Dominum nostrum Jesum Christum, Filium tuum: Qui tecum vivit et
regnat in unitate.

\textbf{English} (Lasance, 1945). Grant, we beseech Thee, O almighty God, that
we who honor the solemnity of Thy blessed martyr Anastasia may experience the
effect of her intercession with Thee. Through our Lord.
\end{massbox}

\begin{massbox}{Christmas Day: Third Mass in the Day (25 December; also Days~5--7 within the Octave)}
\textbf{Latin} (1962 Roman Missal). Concede, quaesumus, omnipotens Deus: ut nos
Unigeniti tui nova per carnem Nativitas liberet; quos sub peccati jugo vetusta
servitus tenet. Per eundem Dominum nostrum Jesum Christum, Filium tuum: Qui
tecum vivit et regnat in unitate Spiritus Sancti, Deus: per omnia saecula
saeculorum. \textbf{R.} Amen.

\textbf{English} (Lasance, 1945). Grant, we beseech Thee, O almighty God, that
the new birth, in the flesh, of Thine only-begotten Son may deliver us whom
slavery from old doth keep under the yoke of sin. Through the same.
\end{massbox}

\begin{massbox}{St.~Stephen, Protomartyr (26 December)}
\textbf{Latin} (1962 Roman Missal). Da nobis, quaesumus, Domine, imitari quod
colimus: ut discamus et inimicos diligere; quia ejus natalicia celebramus, qui
novit etiam pro persecutoribus exorare Dominum nostrum Jesum Christum, Filium
tuum: Qui tecum vivit et regnat in unitate.

\textbf{English} (Lefebvre, 1925). Grant us, we beseech Thee, O Lord, so to
imitate what we revere that we may learn to love even our enemies: for we
celebrate the day of his birth to immortality, who could even plead with our
Lord Jesus Christ Thy Son on behalf of his persecutors. Who with Thee.
\end{massbox}

\begin{massbox}{St.~John, Apostle and Evangelist (27 December)}
\textbf{Latin} (1962 Roman Missal). Ecclesiam tuam, Domine, benignus illustra:
ut, beati Joannis Apostoli tui et Evangelistae illuminata doctrinis, ad dona
perveniat sempiterna. Per Dominum nostrum.

\textbf{English} (Lasance, 1945). Of Thy loving-kindness, O Lord, shed light
upon Thy Church; that, being enlightened by the teachings of blessed John, Thine
apostle and evangelist, it may attain to Thine everlasting gifts. Through our
Lord.
\end{massbox}

\begin{massbox}{The Holy Innocents, Martyrs (28 December)}
\textbf{Latin} (1962 Roman Missal). Deus, cujus hodierna die praeconium
Innocentes Martyres non loquendo, sed moriendo confessi sunt: omnia in nobis
vitiorum mala mortifica; ut fidem tuam, quam lingua nostra loquitur, etiam
moribus vita fateatur. Per Dominum nostrum.

\textbf{English} (Lasance, 1945). O God, Whose praise the martyred innocents did
this day proclaim, not by speaking, but by dying, do to death in us all the
malice of sinfulness, that our lives may also proclaim Thy faith, which our
tongues profess. Through our Lord.
\end{massbox}

\begin{massbox}{Sunday within the Octave of the Nativity (26--31 December)}
\textbf{Latin} (1962 Roman Missal). Omnipotens sempiterne Deus, dirige actus
nostros in beneplacito tuo: ut in nomine dilecti Filii tui mereamur bonis
operibus abundare: Qui tecum vivit et regnat in unitate.

\textbf{English} (Lasance, 1945). O almighty and eternal God, direct our actions
in conformity with Thy good pleasure, that in the name of Thy beloved Son we may
be worthy to abound in good works. Who with Thee.
\end{massbox}

\begin{massbox}{St.~Thomas of Canterbury, Bishop and Martyr (29 December, Commemoration)}
\textbf{Latin} (1962 Roman Missal). Deus, pro cujus Ecclesia gloriosus Pontifex
Thomas gladiis impiorum occubuit: praesta, quaesumus; ut omnes, qui ejus
implorant auxilium, petitionis suae salutarem consequantur effectum. Per Dominum
nostrum Jesum Christum, Filium tuum: Qui tecum vivit et regnat in unitate.

\textbf{English} (Lasance, 1945). O God, for Whose Church the glorious Bishop
Thomas fell by the swords of wicked men, grant, we beseech Thee, that all who
implore his help may obtain the effect of their petition leading to salvation.
Through our Lord.
\end{massbox}

\begin{massbox}{St.~Sylvester~I (31 December) and St.~Telesphorus (5 January), Commemorations}
\textbf{Latin} (1962 Roman Missal, Common of Supreme Pontiffs). Gregem tuum,
Pastor aeterne, placatus intende: et per beatum Silvestrum Summum Pontificem
[for 5~January: \textit{Telesphorum Martyrem tuum atque Summum Pontificem}]
perpetua protectione custodi; quem totius Ecclesiae praestitisti esse pastorem.
Per Dominum nostrum.

\textbf{English} (Lasance, 1945, Common of Supreme Pontiffs). Eternal Shepherd,
do Thou look favorably upon Thy flock, which we beseech Thee to guard and keep
for evermore through the Blessed [Sylvester, / Telesphorus, Thy Martyr and]
Supreme Pontiff, whom Thou didst choose to be the chief shepherd of the whole
Church. Through our Lord.
\end{massbox}

\begin{massbox}{Octave Day of the Nativity (1 January; also Ferias and Our Lady on Saturday, 2--5 January)}
\textbf{Latin} (1962 Roman Missal). Deus, qui salutis aeternae, beatae Mariae
virginitate fecunda, humano generi praemia praestitisti: tribue, quaesumus; ut
ipsam pro nobis intercedere sentiamus, per quam meruimus auctorem vitae
suscipere, Dominum nostrum Jesum Christum, Filium tuum: Qui tecum vivit et
regnat in unitate.

\textbf{English} (Bute, 1908, p.~315). O GOD, Who, by the fruitful virginity of
the Blessed Mary, hast given unto mankind the rewards of everlasting life;
grant, we beseech Thee, that we may continually feel the might of her
intercession, through whom we have worthily received the Author of our life,
our Lord JESUS Christ, Thy Son, Who liveth and reigneth with Thee, in the unity
of the Holy Ghost, one God, world without end. Amen.
\end{massbox}

\begin{massbox}{Feast of the Most Holy Name of Jesus (Sunday 2--5 January, or 2 January)}
\textbf{Latin} (1962 Roman Missal). Deus, qui unigenitum Filium tuum
constituisti humani generis Salvatorem, et Jesum vocari jussisti: concede
propitius; ut, cujus sanctum nomen veneramur in terris, ejus quoque aspectu
perfruamur in caelis. Per eundem Dominum.

\textbf{English} (Bute, 1908). O GOD, Who hast appointed Thine Only-begotten
Son to be the Saviour of mankind, and hast commanded that His Name should be
called JESUS, mercifully grant that we who here on earth do worship that His
most Holy Name may be made glad in heaven by His Presence. Through the Same our
Lord Jesus Christ, Thy Son, Who liveth and reigneth with Thee, in the unity of
the Holy Ghost, One God, world without end. Amen.
\end{massbox}

\subsection{Marian Antiphon, Responsories, and Octave Antiphons}

Latin from the Ratisbon winter Breviary (Pustet, 1888); English from the
Marquess of Bute's \textit{Roman Breviary} (Winter, 1908).

\begin{massbox}{Marian Antiphon of Christmastide: \textit{Alma Redemptoris Mater}}
\textbf{Latin} (Pustet, 1888, pp.~130--131; from First Vespers of Christmas
through Compline of 2~February). Alma Redemptoris Mater, quae pervia coeli Porta
manes, et stella maris, succurre cadenti, Surgere qui curat, populo: tu quae
genuisti, Natura mirante, tuum sanctum Genitorem, Virgo prius ac posterius,
Gabrielis ab ore Sumens illud Ave, peccatorum miserere.

\textbf{From First Vespers of Christmas onward:} \textbf{V.} Post partum Virgo
inviolata permansisti. \textbf{R.} Dei Genitrix intercede pro nobis.
\textit{Oremus.} Deus, qui salutis aeternae beatae Mariae virginitate foecunda,
humano generi praemia praestitisti: tribue, quaesumus; ut ipsam pro nobis
intercedere sentiamus, per quam meruimus auctorem vitae suscipere Dominum
nostrum Jesum Christum Filium tuum. \textbf{R.} Amen.

\textbf{English} (Bute, 1908, pp.~210--211). Maiden! Mother of Him Who redeemed
us, thou that abidest Heaven's open gate, and the Star of the Sea, come, succour
the fallen! Fallen indeed we are, but fain would rise by thy succour. Thou that
beyond nature's course hast borne in time the Eternal; Thou that a Virgin
before, and after that childbirth remainest, From the Archangel's lips the
quickening message receiving, Mother of JESUS and us, turn thine eyes of mercy on
sinners.

\textbf{From First Vespers of Christmas onward:} \textbf{Verse.} After thy
delivery thou still remainest a Virgin undefiled. \textbf{Answer.} Mother of
God, pray for us. \textit{Let us pray.} O GOD, Who, by the fruitful virginity of
the Blessed Mary, hast given unto mankind the rewards of everlasting life;
grant, we beseech Thee, that we may continually feel the might of her
intercession; through whom we have worthily received the Author of our life, our
Lord JESUS Christ Thy Son. \textbf{Answer.} Amen.
\end{massbox}

\begin{massbox}{Christmas Matins Fourth Responsory: \textit{O magnum mysterium}}
\textbf{Latin} (Pustet, 1888, p.~183). \textbf{R.} O magnum mysterium, et
admirabile sacramentum, ut animalia viderent Dominum natum, jacentem in
praesepio: * Beata Virgo, cujus viscera meruerunt portare Dominum Christum.
\textbf{V.} Ave Maria, gratia plena: Dominus tecum. Beata.

\textbf{English} (Bute, 1908, p.~274). How great is this mystery, how wonderful
is the teaching of the faith! The beasts saw the new-born Lord lying in a
manger. Blessed is that Virgin whose womb was made meet to bear our Lord Christ.
\textbf{Verse.} Hail, Mary, full of grace: the Lord is with thee.
\textbf{Answer.} Blessed is that Virgin whose womb was made meet to bear our
Lord Christ.
\end{massbox}

\begin{massbox}{Christmas Second Vespers Magnificat Antiphon: \textit{Hodie Christus natus est}}
\textbf{Latin} (Pustet, 1888, p.~191). Hodie Christus natus est: hodie Salvator
apparuit: hodie in terra canunt Angeli, laetantur Archangeli: hodie exsultant
justi, dicentes: Gloria in excelsis Deo, alleluja.

\textbf{English} (Bute, 1908, pp.~280--281). This day the Christ is born: this
day the Saviour is appeared: * this day the Angels sing praise in the earth and
the Archangels rejoice: this day the righteous are glad and say: Glory to God in
the highest. Alleluia.
\end{massbox}

\begin{massbox}{Octave Day of the Nativity (1 January): The Five \textit{O admirabile commercium} Antiphons}
\textbf{1. Latin} (Pustet, 1888, p.~231). O admirabile commercium: * Creator
generis humani, animatum corpus sumens, de Virgine nasci dignatus est: et
procedens homo sine semine, largitus est nobis suam Deitatem.\\
\textbf{English} (Bute, 1908, p.~315). How wonderful is the dealing here! * The
Maker of man hath taken to Himself a living body, and hath been pleased to be
born of the Virgin: He is come forth conceived without seed, and hath made us
partakers of His Divine nature.\\[0.35em]
\textbf{2. Latin.} Quando natus es * ineffabiliter ex Virgine, tunc impletae
sunt Scripturae: sicut pluvia in vellus descendisti, ut salvum faceres genus
humanum: te laudamus Deus noster.\\
\textbf{English.} When, in an unspeakable manner, * Thou wast born of a Virgin,
the Scriptures were fulfilled; to save man, Thou camest down as dew upon the
fleece: we praise Thee, O our God!\\[0.35em]
\textbf{3. Latin.} Rubum quem viderat Moyses * incombustum, conservatam
agnovimus tuam laudabilem virginitatem: Dei Genitrix, intercede pro nobis.\\
\textbf{English.} Mother of God, when Moses * saw the bush unconsumed, we own
that it was a figure of the preservation of thy most wonderful virginity: pray
for us.\\[0.35em]
\textbf{4. Latin.} Germinavit radix Jesse, * orta est stella ex Jacob: Virgo
peperit Salvatorem: te laudamus Deus noster.\\
\textbf{English.} The Root of Jesse hath blossomed: * the star is arisen out of
Jacob, the Virgin hath given birth to the Saviour: we praise Thee, O our God!\\[0.35em]
\textbf{5. Latin.} Ecce Maria genuit * nobis Salvatorem, quem Joannes videns
exclamavit, dicens: Ecce Agnus Dei, ecce qui tollit peccata mundi, alleluja.\\
\textbf{English.} Behold, Mary hath borne * for us a Saviour; John beheld Him
and cried, saying: Behold, the Lamb of God, behold Him that taketh away the sins
of the world. Alleluia.
\end{massbox}

\subsection{Seasonal Hymns of the Divine Office}

\begin{massbox}{Christmas Vespers and Matins Hymn: \textit{Jesu Redemptor omnium} (Latin)}
\textbf{Latin} (Pustet, 1888, pp.~177--178).\\
Jesu Redemptor omnium,\\
Quem lucis ante originem\\
Parem paternae gloriae\\
Pater supremus edidit.\\[0.4em]
Tu lumen, et splendor Patris,\\
Tu spes perennis omnium,\\
Intende quas fundunt preces\\
Tui per orbem servuli.\\[0.4em]
Memento, rerum Conditor,\\
Nostri quod olim corporis,\\
Sacrata ab alvo Virginis\\
Nascendo, formam sumpseris.\\[0.4em]
Testatur hoc praesens dies,\\
Currens per anni circulum,\\
Quod solus e sinu Patris\\
Mundi salus adveneris.\\[0.4em]
Hunc astra, tellus, aequora,\\
Hunc omne quod coelo subest,\\
Salutis Auctorem novae\\
Novo salutat cantico.\\[0.4em]
Et nos, beata quos sacri\\
Rigavit unda sanguinis;\\
Natalis ob diem tui\\
Hymni tributum solvimus.\\[0.4em]
Jesu, tibi sit gloria,\\
Qui natus es de Virgine,\\
Cum Patre, et almo Spiritu,\\
In sempiterna saecula. Amen.
\end{massbox}

\begin{massbox}{Christmas Vespers and Matins Hymn: \textit{Jesu Redemptor omnium} (English)}
\textbf{English} (Bute, 1908, p.~271).\\
JESUS, the Ransomer of man,\\
Who, ere created light began,\\
Didst from the Sovereign Father spring,\\
His power and glory equalling.\\[0.4em]
Thou brightness of Thy Father's rays,\\
Thou hope and end of all our ways:\\
With gracious ears the prayers attend,\\
Which round the world to Thee ascend.\\[0.4em]
Remember, Lord, that heretofore,\\
When Thee Thy Virgin Mother bore,\\
Thou from her womb didst breathe our air,\\
And human nature for us bear.\\[0.4em]
To Thee, this present solemn day,\\
We yearly adorations pay;\\
The world's Redeemer Thee we own,\\
Descending from Thy Father's throne.\\[0.4em]
The joyful heavens, earth and main,\\
With whatsoever they contain,\\
In new, harmonious accents sing\\
New life restored by the new-born King.\\[0.4em]
We, ransomed by that bloody tide,\\
That issued from Thy sacred side,\\
With double hymns of heart and voice,\\
For this Thy natal day rejoice.\\[0.4em]
To Jesus, from a Virgin sprung,\\
Be glory given, and praises sung:\\
The like to God the Father be,\\
And Holy Ghost eternally. Amen.
\end{massbox}

\begin{massbox}{Christmas Lauds Hymn: \textit{A solis ortus cardine} (Latin)}
\textbf{Latin} (Coelius Sedulius; Pustet, 1888, p.~189).\\
A solis ortus cardine\\
Ad usque terrae limitem,\\
Christum canamus Principem,\\
Natum Maria Virgine.\\[0.4em]
Beatus auctor saeculi\\
Servile corpus induit:\\
Ut carne carnem liberans,\\
Ne perderet quos condidit.\\[0.4em]
Castae Parentis viscera\\
Coelestis intrat gratia:\\
Venter Puellae bajulat\\
Secreta, quae non noverat.\\[0.4em]
Domus pudici pectoris\\
Templum repente fit Dei:\\
Intacta nesciens virum,\\
Concepit alvo Filium.\\[0.4em]
Enititur puerpera,\\
Quem Gabriel praedixerat,\\
Quem ventre Matris gestiens,\\
Baptista clausum senserat.\\[0.4em]
Foeno jacere pertulit:\\
Praesepe non abhorruit:\\
Et lacte modico pastus est,\\
Per quem nec ales esurit.\\[0.4em]
Gaudet chorus coelestium,\\
Et Angeli canunt Deo;\\
Palamque fit pastoribus\\
Pastor, Creator omnium.\\[0.4em]
Jesu, tibi sit gloria,\\
Qui natus es de Virgine,\\
Cum Patre, et almo Spiritu,\\
In sempiterna saecula. Amen.
\end{massbox}

\begin{massbox}{Christmas Lauds Hymn: \textit{A solis ortus cardine} (English)}
\textbf{English} (R.~F.~Littledale, in Bute, 1908, pp.~278--279).\\
From where the sunrise had its birth,\\
Unto the farthest bounds of earth,\\
The Virgin Mary's Child we sing,\\
With notes of gladness, Christ the King.\\[0.4em]
Blest Maker of the world, He came\\
And clad Him in a servant's frame,\\
Our flesh with His own Flesh to aid,\\
Lest those should perish whom He made.\\[0.4em]
His Mother pure becomes the place\\
Wherein abides celestial grace,\\
She bears within her maiden breast\\
A secret by herself unguessed.\\[0.4em]
Soon rises in that modest shrine\\
The Temple of the Lord Divine:\\
The stainless and unwedded one\\
Within her womb conceived the Son.\\[0.4em]
Him in that wondrous birth she bore\\
Whom Gabriel announced before;\\
Whom John while yet unborn perceived,\\
And, leaping in the womb, believed.\\[0.4em]
On hay reclined, the Lord Most High,\\
Within a manger deigned to lie:\\
And He who feeds the birds of air\\
Vouchsafed a little milk to share.\\[0.4em]
The Heavenly choirs now rejoice,\\
The Angels lift to God their voice,\\
The Shepherd Who the world hath made\\
Before the shepherds lies displayed.\\[0.4em]
O Lord, from spotless Virgin sprung,\\
All glory unto Thee be sung,\\
To Father, and to Spirit be\\
Like honour paid eternally. Amen.
\end{massbox}

\begin{massbox}{Holy Innocents Lauds Hymn: \textit{Salvete flores Martyrum}}
\textbf{Latin} (Prudentius; Pustet, 1888, p.~207).\\
Salvete flores Martyrum,\\
Quos lucis ipso in limine\\
Christi insecutor sustulit,\\
Ceu turbo nascentes rosas.\\[0.4em]
Vos prima Christi victima,\\
Grex immolatorum tener,\\
Aram sub ipsam simplices\\
Palma et coronis luditis.\\[0.4em]
Jesu, tibi sit gloria,\\
Qui natus es de Virgine,\\
Cum Patre et almo Spiritu,\\
In sempiterna saecula. Amen.\\[0.5em]
\textbf{English} (J.~M.~Neale, in Bute, 1908, pp.~296--297).\\
All hail! ye infant Martyr flowers!\\
Cut off in life's first dawning hours,\\
As rose-buds snapped in tempest strife,\\
When Herod sought your Saviour's life.\\[0.4em]
You, tender flock of lambs, we sing,\\
First victims slain for Christ your King:\\
Beneath the altar's heavenly ray,\\
With martyr palms and crowns ye play.\\[0.4em]
For their Redemption, glory be,\\
O Jesus, Virgin-born, to Thee!\\
With Father, and with Holy Ghost,\\
For ever from the Martyr-host! Amen.
\end{massbox}

\begin{massbox}{Most Holy Name Vespers Hymn: \textit{Jesu dulcis memoria} (Latin)}
\textbf{Latin} (Pustet, 1888, p.~413).\\
Jesu dulcis memoria,\\
Dans vera cordis gaudia:\\
Sed super mel, et omnia,\\
Ejus dulcis praesentia.\\[0.4em]
Nil canitur suavius,\\
Nil auditur jucundius,\\
Nil cogitatur dulcius,\\
Quam Jesus Dei Filius.\\[0.4em]
Jesu spes poenitentibus,\\
Quam pius es petentibus!\\
Quam bonus te quaerentibus!\\
Sed quid invenientibus?\\[0.4em]
Nec lingua valet dicere,\\
Nec littera exprimere:\\
Expertus potest credere,\\
Quid sit Jesum diligere.\\[0.4em]
Sis Jesu nostrum gaudium,\\
Qui es futurus praemium:\\
Sit nostra in te gloria,\\
Per cuncta semper saecula. Amen.
\end{massbox}

\begin{massbox}{Most Holy Name Vespers Hymn: \textit{Jesu dulcis memoria} (English)}
\textbf{English} (Edward Caswall, in Bute, 1908, p.~365).\\
Jesus! the very thought of Thee\\
With sweetness fills the breast!\\
Yet sweeter far Thy Face to see\\
And in Thy Presence rest.\\[0.4em]
No voice can sing, no heart can frame,\\
Nor can the memory find\\
A sweeter sound than Jesu's Name,\\
The Saviour of mankind!\\[0.4em]
O hope of every contrite heart!\\
O joy of all the meek!\\
To those who fall, how kind Thou art!\\
How good to those who seek!\\[0.4em]
But what to those who find? Ah! this\\
Nor tongue nor pen can show---\\
The love of Jesus, what it is,\\
None but His loved-ones know.\\[0.4em]
Jesus! our only joy be Thou,\\
As Thou our prize shalt be;\\
In Thee be all our glory now,\\
And through eternity! Amen.
\end{massbox}

\subsection{Gospel Canticles of the Divine Office}

Latin from the Clementine Vulgate collated against liturgical editions
(restoring the received reading \textit{sicut locutus est} at Luke~1:70; see
Appendix~\ref{app:about}); English from the American Douay-Rheims Bible (1899).
In the Divine Office each canticle closes with the \textit{Gloria Patri}.

\begin{massbox}{\textit{Magnificat} (Vespers, Luke 1:46--55)}
\textbf{Latin.} Magnificat anima mea Dominum: et exsultavit spiritus meus in Deo
salutari meo. Quia respexit humilitatem ancill\ae{} su\ae{}: ecce enim ex hoc
beatam me dicent omnes generationes, quia fecit mihi magna qui potens est: et
sanctum nomen ejus, et misericordia ejus a progenie in progenies timentibus
eum. Fecit potentiam in brachio suo: dispersit superbos mente cordis sui.
Deposuit potentes de sede, et exaltavit humiles. Esurientes implevit bonis: et
divites dimisit inanes. Suscepit Isra\"el puerum suum, recordatus
misericordi\ae{} su\ae{}: sicut locutus est ad patres nostros, Abraham et semini
ejus in s\ae cula.

\textbf{English.} My soul doth magnify the Lord. And my spirit hath rejoiced in
God my Saviour. Because he hath regarded the humility of his handmaid; for
behold from henceforth all generations shall call me blessed. Because he that is
mighty, hath done great things to me; and holy is his name. And his mercy is
from generation unto generations, to them that fear him. He hath shewed might in
his arm: he hath scattered the proud in the conceit of their heart. He hath put
down the mighty from their seat, and hath exalted the humble. He hath filled the
hungry with good things; and the rich he hath sent empty away. He hath received
Israel his servant, being mindful of his mercy: As he spoke to our fathers, to
Abraham and to his seed for ever.
\end{massbox}

\begin{massbox}{\textit{Benedictus} (Lauds, Luke 1:68--79)}
\textbf{Latin.} Benedictus Dominus Deus Isra\"el, quia visitavit, et fecit
redemptionem plebis su\ae{}: et erexit cornu salutis nobis in domo David pueri
sui, sicut locutus est per os sanctorum, qui a s\ae culo sunt, prophetarum ejus:
salutem ex inimicis nostris, et de manu omnium qui oderunt nos: ad faciendam
misericordiam cum patribus nostris: et memorari testamenti sui sancti:
jusjurandum, quod juravit ad Abraham patrem nostrum, daturum se nobis ut sine
timore, de manu inimicorum nostrorum liberati, serviamus illi in sanctitate et
justitia coram ipso, omnibus diebus nostris. Et tu puer, propheta Altissimi
vocaberis: pr\ae ibis enim ante faciem Domini parare vias ejus, ad dandam
scientiam salutis plebi ejus in remissionem peccatorum eorum per viscera
misericordi\ae{} Dei nostri, in quibus visitavit nos, oriens ex alto: illuminare
his qui in tenebris et in umbra mortis sedent: ad dirigendos pedes nostros in
viam pacis.

\textbf{English.} Blessed be the Lord God of Israel; because he hath visited and
wrought the redemption of his people: And hath raised up an horn of salvation to
us, in the house of David his servant: As he spoke by the mouth of his holy
prophets, who are from the beginning: Salvation from our enemies, and from the
hand of all that hate us: To perform mercy to our fathers, and to remember his
holy testament, The oath, which he swore to Abraham our father, that he would
grant to us, That being delivered from the hand of our enemies, we may serve him
without fear, In holiness and justice before him, all our days. And thou, child,
shalt be called the prophet of the Highest: for thou shalt go before the face of
the Lord to prepare his ways: To give knowledge of salvation to his people, unto
the remission of their sins: Through the bowels of the mercy of our God, in
which the Orient from on high hath visited us: To enlighten them that sit in
darkness, and in the shadow of death: to direct our feet into the way of peace.
\end{massbox}

\begin{massbox}{\textit{Nunc Dimittis} (Compline, Luke 2:29--32)}
\textbf{Latin.} Nunc dimittis servum tuum Domine, secundum verbum tuum in pace:
quia viderunt oculi mei salutare tuum, quod parasti ante faciem omnium
populorum: lumen ad revelationem gentium, et gloriam plebis tu\ae{} Isra\"el.

\textbf{English.} Now thou dost dismiss thy servant, O Lord, according to thy
word in peace; Because my eyes have seen thy salvation, Which thou hast prepared
before the face of all peoples: A light to the revelation of the Gentiles, and
the glory of thy people Israel.
\end{massbox}
